\begin{papercard}{light2018}
\cf{Stages} \stagetag{D} \stagetag{W} (also \stagetag{S} \stagetag{R}). The founding statement of the \emph{app walkthrough method} in media studies. It has two parts: the \emph{environment of expected use} (vision, operating model, governance) and a \emph{technical walkthrough} through registration, everyday use and leaving. It is a critical, interpretive method, unlike HCI usability walkthroughs.

\cf{Purpose and theory} Engages with an app's interface to examine technological mechanisms and embedded cultural references, i.e.\ how the app guides users and shapes their experience \pc{light2018}{2}{The walkthrough method is a way of engaging directly with an}. Its core is step-by-step observation and documentation of screens, features and flows, which slows down mundane actions to make them available for critical analysis \pc{light2018}{2}{The core of this method involves the step-by-step observation and documentation}. It is explicitly \emph{not} a test of whether users respond as designers intended. It illuminates the material traces of those intentions \pc{light2018}{6}{The walkthrough method we propose is used, not to test whether users respond to an interface}. Grounded in Actor-Network Theory: interface elements are non-human actors, and \emph{mediators} transform meaning \pc{light2018}{6}{Intermediaries pass meaning along unchanged throughout a network of relations}. Affordances are examined at several scales, from buttons to operating system and other apps \pc{light2018}{6}{The walkthrough examines affordances at multiple levels of scale}. Cultural studies (``circuit of culture'') adds representation, identity, production, consumption and regulation \pc{light2018}{8}{how it is produced and consumed, and what mechanisms regulate its distribution and use}.

\cf{Development and evidence} Developed in studies of dating and hook-up apps, then refined in methods workshops with more than 250 academics and graduate students across domains including health \pc{light2018}{4}{a series of research methods workshops involving more than 250 academics}. It is adaptable in whole or in part \pc{light2018}{4}{adaptable – in whole or in part – to a range of research questions}. The opening example is a menstrual tracker (Clue) whose sex icons reduce sexual activity to binaries and assume heterosexuality \pc{light2018}{2}{The icons also reinforce cultural norms of heterosexuality}.

\cf{Part 1: environment of expected use} Context examined before or alongside the walkthrough \pc{light2018}{3}{This work to establish the envi}. \emph{Vision}: purpose, target users and scenarios of use, from app-store descriptions, websites, blogs, marketing and press releases \pc{light2018}{9}{An app’s vision involves its purpose, target user base and scenarios of use}. \emph{Operating model}: business strategy and revenue, including personal data exchanged for services. Sources include store price, in-app purchase menus, recruitment and press kits, LinkedIn, Crunchbase and SEC filings \pc{light2018}{10}{Since many apps allow access to services in exchange for personal data that can be sold to advertisers}, \pc{light2018}{10}{The app’s price in app stores, in-app purchase menus}. \emph{Governance}: how the provider manages user activity through rules, hard-coded features, TOS and FAQ; study TOS over time with the Wayback Machine \pc{light2018}{11}{conceptions of governance, especially when studied over time using tools like the Internet Archive}. The platform and regional legal status also govern use \pc{light2018}{11}{the platform through which the app is delivered may act as a governance mechanism}.

\cf{Part 2: technical walkthrough} The central data-gathering procedure: the researcher takes a user's position while applying an analytical eye to acquiring, registering, using and discontinuing \pc{light2018}{11}{Walking through the app requires the researcher assume a user’s position}. Data: detailed field notes, screenshots, screen video and audio recordings of one's own thoughts \pc{light2018}{11}{producing a collection of data by generating detailed field notes and recordings}. Four \textbf{mediator characteristics} guide attention. (1) \emph{User interface arrangement}, e.g.\ smaller, harder-to-find ``Report'' or ``Share'' buttons \pc{light2018}{11}{User interface arrangement. How the app guides users through activities}. (2) \emph{Functions and features}: pop-ups, compulsory fields, account-linking requests \pc{light2018}{11}{Groups of arrangements that mandate or enable an activ}. (3) \emph{Textual content and tone}: embedded text, the order of menu options, available categories \pc{light2018}{11}{Textual content and tone. More than instructions}. (4) \emph{Symbolic representation}: look and feel, branding, colour, font \pc{light2018}{12}{A semiotic approach to examining the look and feel of the app}.

\cf{Three stages of use} \emph{Registration and entry}: routes (web or mobile), third-party login such as Facebook, and requests for contacts, photos and location \pc{light2018}{12}{Other examples of common registration steps include requests for access}. Vision, operating model and governance are often salient here, e.g.\ ``I agree'' terms screens \pc{light2018}{12}{Registration is also a stage where an app generally com}. \emph{Everyday use}: record functionality, options and affordances. This may need multiple profiles, and the \emph{flow} (order of screens) is data too \pc{light2018}{13}{Paying attention not only to the app’s features but also to the flow of activity}. The menu can anchor exploration, or the focus can be on the most common activities \pc{light2018}{14}{the researcher may use the menu to anchor exploration}. Static screens miss mediation, so flows such as location permission requests must be followed \pc{light2018}{14}{Since attending only to static interface elements risks missing important points of mediation}. \emph{Suspension, closure and leaving}: logging out, hiding profiles, data removal. Example: deleting the Facebook app does not delete server data, and deletion flows use affective persuasion \pc{light2018}{14}{this action deletes data from the user’s device but not from Facebook’s database}. Leaving can reveal how apps retain value from users even after they leave \pc{light2018}{15}{retain value from users even after they leave}.

\cf{Unexpected practices} Without user data, the method can still reveal user-led appropriation (third-party apps, hacks, critique), sometimes requiring walkthroughs of third-party apps \pc{light2018}{15}{It may be necessary to perform the walkthrough with third-party apps}.

\cf{Ethics and limitations} Dummy accounts may be mistaken for real users, so plan for interaction attempts in advance \pc{light2018}{16}{Strategies for managing potential user interactions need to be considered before any empirical work}. Other users' data seen during the walkthrough should be anonymised as in observation research \pc{light2018}{16}{User data should be treated similar to data gathered from observation methods}. Limits: no direct user data; some features may be programming by-products. Combine with reviews, interviews or API data. Because apps are not stable artefacts, the walkthrough may need repeating across updates \pc{light2018}{16}{Since apps are not stabilised artefacts, it may be necessary to conduct the walkthrough multiple times}.

\cf{Use for this thesis} The canonical source for an app-centred walkthrough and for \emph{stage D}. The three stages (registration, everyday use, leaving) map directly onto privacy tasks: consent at onboarding, settings during use, deletion and export at exit. The environment of expected use gives a structured capture of store listing, website, policy and business model, with the Wayback Machine for versions. Record screenshots, screen video and field notes, and anonymise any third-party content. Light et al.\ is interpretive, so pairing it with the CW questions (Wharton) or with a codebook is a design choice the thesis must justify.
\end{papercard}
